% !TeX spellcheck = en-US
% !TeX encoding = utf8
% !TeX program = pdflatex
% !BIB program = biber
% -*- coding:utf-8 mod:LaTeX -*-


% vv  scroll down to line 200 for content  vv


\let\ifdeutsch\iffalse
\let\ifenglisch\iftrue
\input{pre-documentclass}
\documentclass[
  % fontsize=11pt is the standard
  a4paper,  % Standard format - only KOMAScript uses paper=a4 - https://tex.stackexchange.com/a/61044/9075
  twoside,  % we are optimizing for both screen and two-sided printing. So the page numbers will jump, but the content is configured to stay in the middle (by using the geometry package)
  bibliography=totoc,
  %               idxtotoc,   %Index ins Inhaltsverzeichnis
  %               liststotoc, %List of X ins Inhaltsverzeichnis, mit liststotocnumbered werden die Abbildungsverzeichnisse nummeriert
  headsepline,
  cleardoublepage=empty,
  parskip=half,
  %               draft    % um zu sehen, wo noch nachgebessert werden muss - wichtig, da Bindungskorrektur mit drin
  draft=false
]{scrbook}
\input{config}


\usepackage[
  title={Towards Human-Aligned Chess Engines: Replicating Human Player Behavior across Different Skill Levels},
  author={Andrei Todiraș},
  type=bachelor,
  institute=Institute\ for\ Artificial\ Intelligence\\Department\ Machine\ Learning\ for\ Simulation\ Sciences, % or other institute names - or just a plain string using {Demo\\Demo...}
  course={Informatik},
  examiner={Prof.\ Dr.\ Mathias Niepert},
  supervisor={Andrei Manolache,\ M.Sc.},
  startdate={01.01.2026},
  enddate={31.03.2026}
]{scientific-thesis-cover}

\input{acronyms}

\makeindex

\begin{document}

%tex4ht-Konvertierung verschönern
\iftex4ht
  % tell tex4ht to create pictures also for formulas starting with '$'
  % WARNING: a tex4ht run now takes forever!
  \Configure{$}{\PicMath}{\EndPicMath}{}
  %$ % <- syntax highlighting fix for emacs
  \Css{body {text-align:justify;}}

  %conversion of .pdf to .png
  \Configure{graphics*}
  {pdf}
  {\Needs{"convert \csname Gin@base\endcsname.pdf
      \csname Gin@base\endcsname.png"}%
    \Picture[pict]{\csname Gin@base\endcsname.png}%
  }
\fi

%\VerbatimFootnotes %verbatim text in Fußnoten erlauben. Geht normalerweise nicht.

\input{commands}
\pagenumbering{arabic}
\Titelblatt

%Eigener Seitenstil fuer die Kurzfassung und das Inhaltsverzeichnis
\deftriplepagestyle{preamble}{}{}{}{}{}{\pagemark}
%Doku zu deftriplepagestyle: scrguide.pdf
\pagestyle{preamble}
\renewcommand*{\chapterpagestyle}{preamble}



%Kurzfassung / abstract
%auch im Stil vom Inhaltsverzeichnis
\section*{Abstract}
Human-aligned chess engines are designed to replicate human player behavior across different skill levels, providing more relatable and understandable chess analysis and recommendations.
As the field is still in its early stages, there are many open questions regarding the best methods for achieving this goal, as well as the challenges and limitations of these methods.
The goal of this thesis is to investigate the potential of two methods for improving human-aligned chess engines, unused in the current state-of-the-art models: behavior stylometry and adversarial training.

Behavior stylometry filters the training data to only include games played in a certain style, thus allowing the model to learn from a more homogeneous dataset and creating a more nuanced engine.
Adversarial training, on the other hand, modifies the training process by introducing a discriminator that tries to distinguish between human and engine-generated moves, thus aiming to improve the generator's predictions.
This thesis presents the whole process of implementing these two methods, discussing not just the final results, but also the challenges and insights gained along the way.

The results indicate that while behavior stylometry does not lead to a significant improvement in the model's accuracy, it does modify the model's predictions in a way that is consistent with the intended style.
Adversarial training proves to be a much more difficult task, not only in its implementation, but also in result predictability and explainability.
While the final model yields a discriminator that can distinguish between human and engine-generated moves with high accuracy, the generator's predictions tend to be blunder-prone, degrading the model's performance at higher Elo levels.

We discuss the reasons behind these outcomes, alternative solutions, and key insights, advocating psychologically-informed evaluation metrics and training methods as a promising direction for future research in the field of human-aligned engines.
\clearpage
%Solely for German courses of study
\section*{Kurzfassung}
Menschlich-ausgerichtete Schachengines sollen das Verhalten menschlicher Spieler auf verschiedenen Spielstärken nachbilden und so eine anschaulichere und verständlichere Schachanalyse sowie Empfehlungen liefern.
Da sich dieses Forschungsgebiet noch in einem frühen Stadium befindet, gibt es viele offene Fragen hinsichtlich der besten Methoden zur Erreichung dieses Ziels sowie der Herausforderungen und Grenzen dieser Methoden.
Das Ziel dieser Arbeit ist es, das Potenzial zweier Methoden zur Verbesserung menschlich-ausgerichteter Schachengines zu untersuchen, die in den aktuellen state-of-the-art-Modellen noch nicht zum Einsatz kommen: Verhaltensstylometrie und adversarielles Training.

Die Verhaltensstylometrie filtert die Trainingsdaten so, dass nur Partien eines bestimmten Stils berücksichtigt werden. Dadurch kann das Modell aus einem homogeneren Datensatz lernen und eine differenziertere Engine entwickeln.
Das adversarielles Training hingegen modifiziert den Trainingsprozess durch die Einführung eines Diskriminators, der versucht, zwischen menschlichen und Engine-generierten Zügen zu unterscheiden, und zielt somit darauf ab, die Vorhersagen des Generators zu verbessern.
Diese Arbeit stellt den gesamten Prozess der Implementierung dieser beiden Methoden vor und diskutiert nicht nur die Endergebnisse, sondern auch die Herausforderungen und Erkenntnisse, die dabei gewonnen wurden.

Die Ergebnisse zeigen, dass die Verhaltensstylometrie zwar nicht zu einer signifikanten Verbesserung der Modellgenauigkeit führt, die Vorhersagen des Modells jedoch in einer Weise modifiziert, die mit dem beabsichtigten Stil übereinstimmt.
Das adversarielle Training erweist sich als eine wesentlich schwierigere Aufgabe, nicht nur hinsichtlich der Implementierung, sondern auch in Bezug auf die Vorhersagbarkeit und Erklärbarkeit der Ergebnisse.
Während das endgültige Modell einen Diskriminator liefert, der mit hoher Genauigkeit zwischen menschlichen und Engine-generierten Zügen unterscheiden kann, neigen die Vorhersagen des Generators zu Fehlern, was die Leistung des Modells auf höheren Elo-Niveaus beeinträchtigt.

Wir diskutieren die Gründe für diese Ergebnisse, alternative Lösungen und wichtige Erkenntnisse und plädieren für psychologisch fundierte Bewertungsmetriken und Trainingsmethoden als vielversprechende Richtung für die zukünftige Forschung im Bereich der menschenorientierten Engines.

%\cleardoublepage


% BEGIN: Verzeichnisse

\iftex4ht
\else
  \microtypesetup{protrusion=false}
\fi

%%%
% Literaturverzeichnis ins TOC mit aufnehmen, aber nur wenn nichts anderes mehr hilft!
% \addcontentsline{toc}{chapter}{Literaturverzeichnis}
%
% oder zB
%\addcontentsline{toc}{section}{Abkürzungsverzeichnis}
%
%%%

%Produce table of contents
%
%In case you have trouble with headings reaching into the page numbers, enable the following three lines.
%Hint by http://golatex.de/inhaltsverzeichnis-schreibt-ueber-rand-t3106.html
%
%\makeatletter
%\renewcommand{\@pnumwidth}{2em}
%\makeatother
%
%\tableofcontents

% Bei einem ungünstigen Seitenumbruch im Inhaltsverzeichnis, kann dieser mit
% \addtocontents{toc}{\protect\newpage}
% an der passenden Stelle im Fließtext erzwungen werden.

%\listoffigures
%\listoftables

%Wird nur bei Verwendung von der lstlisting-Umgebung mit dem "caption"-Parameter benoetigt
%\lstlistoflistings
%ansonsten:
%\ifdeutsch
%  \listof{Listing}{Verzeichnis der Listings}
%\else
%  \listof{Listing}{List of Listings}
%\fi

%mittels \newfloat wurde die Algorithmus-Gleitumgebung definiert.
%Mit folgendem Befehl werden alle floats dieses Typs ausgegeben
%\ifdeutsch
%  \listof{Algorithmus}{Verzeichnis der Algorithmen}
%\else
%  \listof{Algorithmus}{List of Algorithms}
%\fi
%\listofalgorithms %Ist nur für Algorithmen, die mittels \begin{algorithm} umschlossen werden, nötig

% Abkürzungsverzeichnis
%\printnoidxglossaries

\iftex4ht
\else
  %Optischen Randausgleich und Grauwertkorrektur wieder aktivieren
  \microtypesetup{protrusion=true}
\fi

% END: Verzeichnisse


% Headline and footline
\renewcommand*{\chapterpagestyle}{scrplain}
\pagestyle{scrheadings}
\pagestyle{scrheadings}
\ihead[]{}
\chead[]{}
\ohead[]{\headmark}
\cfoot[]{}
\ofoot[\usekomafont{pagenumber}\thepage]{\usekomafont{pagenumber}\thepage}
\ifoot[]{}


%% vv  scroll down for content  vv %%































%%%%%%%%%%%%%%%%%%%%%%%%%%%%%%%%%%%%%%%%%%%%%%%%%%%%%%%%%%%%%%%%%%%%%%%%%%%%%%
%
% Main content starts here
%
%%%%%%%%%%%%%%%%%%%%%%%%%%%%%%%%%%%%%%%%%%%%%%%%%%%%%%%%%%%%%%%%%%%%%%%%%%%%%%
\chapter{Introduction}
\label{chapter:introduction}
Chess has been at the forefront of artificial intelligence research for decades, serving not only as a means of presenting the capabilities of AI agents to the wide public, but also as a challenging domain for developing and testing new algorithms.
Even before artificial intelligence as a field was formally established, early computer scientists were fascinated by the idea of creating machines that could play chess.
From Alan Turing's Turochamp in the 1940s to IBM's Deep Blue in the 1990s, the last century saw tremendous advancements in chess-playing computers, though these were more of a demonstration of brute-force computational power rather than true intelligence \cite{chess_computer_history}.
DeepBlue was the culmination of this era, defeating world champion at the time Garry Kasparov in 1997 by employing 32 processors to evaluate 200 million positions per second \cite{deepblue}.
Ever since this legendary match, it has been clear that computers can surpass human performance in chess.
However, the development of chess engines has not stopped there, with modern engines like Stockfish \cite{stockfish} and Leela Chess Zero \cite{leela} now far exceeding the strongest human players, without even needing massive computational resources.

While these engines are incredibly strong and have changed the way chess is played and studied, they have also introduced a new challenge: their playstyle is often very different from that of human players, making it difficult for humans to learn from and interact with them.
Of course, professional players have adapted to this new reality, but for amateur players, playing against such engines yields little educational value, and even interpreting their analysis and recommendations can be challenging due to the great depth and complexity of their calculations.
Considering that the great majority of players today are either amateurs or total beginners (see Figure \ref{fig:chess_player_distribution}), it is clear that the demand for educational content and training tools is high.
Traditional learning material in the form of opening books or endgame studies, however, is losing relevance in a world more and more reliant on artificial intelligence.
As such, there is a need for engines that not only play at a high level but also present their thought process in a way that is understandable and useful for human learning.

\begin{figure}[H]
    \centering
    \includegraphics[width=0.4\textwidth]{chess_com_stats.pdf}
    \caption{Distribution of chess.com users by rating as of March 2026 \cite{chess_com_stats}. Advanced players are considered to have a rating above 1800 (under 10\% of all players)}
    \label{fig:chess_player_distribution}
\end{figure}

Another big current problem in the world of chess is represented by the rise of cheating, especially in online chess.
With the increasing strength and accessibility of chess engines, it has become easier for players to use them to gain an unfair advantage, particularly in online games where there is no physical presence to monitor the players.
Anti-cheating measures have been implemented by online platforms, but they rely heavily on statistical analysis such as winrate, average accuracy compared to engine recommendations, and manual review of suspicious games \cite{chess_com_anticheat}.
An engine that could mimic human play at the level of an amateur player would be a powerful tool for anti-cheating measures, as it could be used to generate synthetic data for training anti-cheating algorithms, or even to directly evaluate the human-likeness of a player's moves in real-time.

To address these challenges, the field of ``human-aligned'' chess engines has emerged, focusing on creating systems that emulate human play.
The field is still in its incipient phase, and this thesis aims to contribute to it by investigating different methods of improving such engines, namely behavior stylometry and adversarial training.
While these methods have been widely studied in other domains, and even applied to chess engines in some cases, to the best of our knowledge they have not been used in human-aligned engines before in the way we propose here.
This thesis thus explores these methods in depth, evaluates their effectiveness, and provides insights into their potential for creating more human-like chess engines.

In order to explain the context of this work, we begin with an analysis of existing chess engines in Chapter \ref{chapter:foundations}.
Then we present the related work, methodology and evaluation of the two methods in Chapters \ref{chapter:behavior_stylometry} and \ref{chapter:adversarial_training} respectively.
To conclude, final discussions and ideas for future work are included in Chapter \ref{chapter:conclusion}.
\chapter{Foundations}
\label{chapter:foundations}

\section{Human vs. Engine Move}
To begin understanding the reason behind the development of human-aligned chess engines, the most important aspect to consider is the difference between human and engine moves.
For people who are not deeply familiar with chess, this difference may not be immediately obvious, as humans and engines often choose the same moves in a given position.
Behind any human move lies a complex thought process, involving general intuition, pattern recognition, and strategic planning.
Humans rely on previous experience, learned principles, and rather abstract concepts to guide their decision-making.
Even outer factors, such as time pressure, opponent rating, must-win situation or not, can influence the choice of move, but we will not consider these in this work.

Focusing on the game alone, the best example to illustrate human abstract thinking is an endgame position.
When one learns endgame principles, they do not memorize specific moves, but rather general positions which are known to be winning or drawing, or abstract ideas and concepts.
One such concept is ``king opposition'' and it refers to placing the king on the same rank or file in order to force the opponent to choose one direction to move their own king, and then go around it the other way, infiltrating.
Figure \ref{fig:king_opposition} shows a concrete example, where it is black to move.
The position comes from a blitz game played between two international masters (the second highest title in chess), so a very high-level game.
\begin{figure}[H]
  \centering
  \includegraphics[width=0.4\textwidth]{king_opposition.png}
  \caption{Ke6 achieves ``king opposition'' and white's king cannot advance. Stockfish recommends h5, allowing Kd5}
  \label{fig:king_opposition}
\end{figure}
The move Ke6 is by far the most natural move, because it gets the king opposition, meaning it stops the white king from advancing anywhere, especially on the d5 square.
This is the move that was played in the game as well.
However, Stockfish says that it is a losing move, instead suggesting h5 as the only way to draw, which is rather unnatural, because any player is taught that in this kind of endgames, the most important thing is to not allow the enemy king forward.

To give a different example, consider the position in Figure \ref{fig:human_vs_engine}, with black to move.
\begin{figure}[H]
    \centering
    \includegraphics[width=0.4\textwidth]{Human_move.png}
    \caption{Human vs. Engine Move: White attacks the f7 pawn, but Stockfish recommends b5}
    \label{fig:human_vs_engine}
\end{figure}
One can see that white has a powerful attack, not only threatening to win the pawn on f7, but also preparing checkmating ideas by pushing the g pawn or bringing the c3 knight to d5.
Stockfish, instead of defending the pawn, recommends a counterattack on the other side of the board with b5, a move that not only gives up material and allows the white queen to get closer to the black king, but also does not create a threat in return.
Maia, a human-aligned engine, chooses a much more ``natural'', human-like move, which is e5, creating the threat of a fork in the middle of the board, while also opening the diagonal for the black queen.
The reason why Stockfish's move is objectively better (in that it leads to a draw instead of a loss) can only be seen 20 moves later, when the attack fizzles out and that advanced pawn becomes an asset in the endgame.
Such depth of calculation and ignoring of immediate threats is simply not human though, and it is hard for a beginner or intermediate player to learn from such moves.
This is the main motivation behind human-aligned engines: to create chess engines that play in a more human-like manner, making moves that are understandable and instructive for human players.

\section{Modern Chess Engines}
In order to understand the challenges and methods involved in creating human-aligned chess engines, it is important to first look at how regular modern engines work.
We refer to these engines as ``modern'' because their development was not defined by hardware limitations, and they do not rely on brute-force techniques and massive amounts of computation to achieve their strength.
Instead, the algorithms that they employ are what make them strong, and they are able to run on consumer-grade hardware.

As will be seen later (Section \ref{section:existing_engines}), human-aligned engines build upon the same foundations as these rather than introducing entirely new concepts, so a solid understanding of their inner workings is essential.
All modern engines share a common architecture, consisting of two components: a neural network that, given a board position, outputs a so-called evaluation function, and a search algorithm that uses this evaluation function to explore the game tree and select the best move.
This architecture was first introduced in 2017 by DeepMind's AlphaZero \cite{alphazero} and has since become the standard for chess engines. 
The neural network is trained on large amounts (tens and even hundreds of millions) of self-generated games, that is, games played between instances of the engine itself, and is responsible for generating general patterns and understanding of the game.
The search algorithm, on the other hand, is responsible for exploring specific variations and calculating concrete move sequences, impacting the performance and strength of the engine.
\begin{figure}[H]
  \centering
  \includegraphics[width=\textwidth]{AlphaZero.pdf}
  \caption{AlphaZero Architecture \cite{alphazero_cheat_sheet}: the NN takes the board position as input and outputs a move probability distribution and a value estimation, which are then used by the MCTS search algorithm to explore the game tree and select the best move.}
  \label{fig:AlphaZero}
\end{figure}

This, however, is where the similarities end.
When it comes to the neural network, Stockfish uses a very shallow network with only a few layers \cite{stockfish}, while Leela Chess Zero employs a deep convolutional neural network with many layers and millions of parameters \cite{leela}.
In general, the CNN-based engines are considered to be more human-like in their playstyle, generating more clear general patterns that are easier for humans to understand and apply in their own games.
The search algorithms also differ significantly, with Stockfish relying on a highly optimized alpha-beta pruning algorithm, while Leela Chess Zero uses a Monte Carlo Tree Search (MCTS) approach.
This is the main reason why Stockfish is the strongest and most popular engine: its alpha-beta search, based on decades of research and optimization, is simply more efficient and effective than MCTS in the context of chess.
However, MCTS has the advantage of being more flexible and adaptable, allowing for more human-like decision-making and exploration of the game tree.
That is why research in human-aligned engines has focused on Leela Chess Zero as a base model, as it is more amenable to modifications that can lead to more human-like behavior.

\section{Existing Human-Aligned Engines}
\label{section:existing_engines}
As far as we have been able to determine, there are four human-aligned chess engines that have been well documented in literature: Maia \cite{maia} and its successor Maia2 \cite{maia2} (which we will be using in this thesis), Allie \cite{allie} and Harmonia \cite{harmonia}.
Maia was the first engine of this type, developed in 2020, and it came as a simple modification of Leela Chess Zero.
The idea was extremely straightforward: instead of letting the neural network train on self-play games, train it on human games.
Since there are huge databases with billions of human games available (a very large number being of course necessary for training neural networks), this wouldn't be a problem.
The more fascinating aspect was the search algorithm: the authors found that better results were obtained by just removing it entirely, and letting the neural network select moves directly.
The cause of this, as the authors of Allie later discovered, is that Maia only aimed at emulating amateur-level players, not professionals.
At the high levels, where players align very closely with Stockfish evaluations, the search algorithm is necessary to capture the small differences in move quality.
Thus Allie was developed, keeping almost the same neural network as Maia, but introducing a time-adaptive MCTS search algorithm that takes into account how long human players at different levels take to make moves. 
Harmonia, on the other hand, changed the neural network architecture itself, using an encoder-only transformer that proved to be much more efficient in size and training time.

We did not focus on Allie or Harmonia in this thesis, the reason being that their codebases are not publicly available, and we would not be able to guarantee a faithful reproduction.
Instead, this work relied on Maia2, developed by the same team as Maia, improving upon it by taking players' rating into account in a different way.
Instead of only using games played by humans of a certain rating range and thus training many separate models, Maia2 introduced rating as an additional input to the neural network, allowing a single model to adapt its playstyle to different levels.
Thus a single engine model is created, which plays differently not just based on its rating (the player level it is emulating), but also based on the opponent's rating, leading to a much more dynamic and human-like behavior.
This is because a human plays differently depending on the opponent: against a weaker opponent, they may play more aggressively and take more risks in order to force wins, while against a stronger opponent, they may play more cautiously and defensively.
\begin{figure}[H]
  \hspace{-2cm}
  \includegraphics[width=1.3\textwidth]{Maia2_Architecture.pdf}
  \caption{Maia2 Architecture: convolutional layers process the board position, which is then combined with the rating input and passed through fully connected attention layers to produce the move probability distribution and value estimation}
  \label{fig:Maia2}
\end{figure}

The results obtained by Maia2 are impressive, with the engine being able to closely mimic human playstyles at different levels, and even outperforming Maia and Allie in terms of move prediction accuracy.
The codebase is publicly available\footnote{\url{https://github.com/CSSLab/maia2}} and represented the base for the methods proposed in this thesis.
\chapter{Behavior Stylometry}
\label{chapter:behavior_stylometry}

\section{Related Work}
\label{section:related_work_behavior_stylometry}
Behavior stylometry is the study of identifying individual characteristics based on patterns of behavior.
Historically, it has been applied in various fields, from psychology to cybersecurity.
In the domain of machine learning, behavior stylometry has been widely employed for many different applications, such as personalized author identification and profiling \cite{author_identification}, user authentication \cite{user_authentication}, or fraud detection \cite{fraud_detection}.
Of course, these are just a few examples, and it is easy to see how there are countless other applications where it can be useful.
In the most general sense, behavior stylometry involves the collection and analysis of specific data that reflects an individual's behavior.
This is why it is often used in conjunction with machine learning techniques, as they can effectively process and analyze large datasets to identify patterns and make predictions.
While traditional stylometry initially focused on text (linguistic features), because textual data is the most readily available and easily processed, the field has since expanded to image recognition \cite{image_recognition} or speech recognition \cite{speech_recognition}.

For chess, there have been two main directions of research.
The first direction focuses on identifying players based on a number of games that they have played \cite{chess_stylometry1}.
This is done by straightforward neural networks of various architectures that take as input a sequence of moves and output the predicted player, after being trained on a large dataset of games.
Manual feature extraction is usually minimal, the network learning the relevant features by itself, relying on the sheer amount of data.
For this reason, these methods were not particularly useful for our purposes, since we wanted more explainability and interpretability of the results, as well as being able to work with a limited number of games.
The second direction of research focuses on creating engines that mimic the playing style of a specific player \cite{chess_stylometry2}.
These engines work with much smaller datasets (only games of the player to be mimicked) and rely more on manual feature extraction, which made them more relevant to this work.

\section{Methodology}
\label{section:methodology_behavior_stylometry}
In this thesis, we focused on the second direction, aiming to apply behavior stylometry techniques to improve human-aligned chess engines.
Starting from the Maia2 engine \cite{maia2}, we trained models on different datasets representing different styles of play, namely aggressive, defensive, and mixed.
The work was based mainly on \cite{chess_stylometry2}, but with some modifications to better suit our needs.
The author of \cite{chess_stylometry2} defined a set of 8 features that were used to split players into 5 categories (each similar to the style of one famous grandmaster).
Instead, we defined only 3 categories: aggressive, defensive, and mixed (playstyle that does not clearly fit in either of the other two), because we wanted to have a more general classification that can be applied to a wider range of players.
In this way, the resulting models would not be tied to a specific player's style, but rather to a more general style, combining behavior stylometry with the advantage of training on large datasets.
As such, ``aggressive'' refers to very tactical games, with dynamic openings, tactical middlegames, attacks on the king, piece sacrifices, and checkmating patterns.
``Defensive'' refers to more positional games, with either many trades and long endgames, or closed positions with very few trades, where strategic maneuvering is more important than direct attacks.
Of course, it is impossible to perfectly distinguish between these styles, and many games have elements of both, which is why the ``mixed'' category was necessary.

To better understand the categories, let's take two examples.
The first game\footnote{\url{https://lichess.org/d3UVZ3Th\#51}} started with a gambit line and a very dynamic position (Figure \ref{fig:game1_opening}), with multiple pins and tactics available to both sides.
In the middlegame, white blundered, missing that black's capture on f3 comes with a check and thus their pin on the e8 rook does not actually work (Figure \ref{fig:game1_middle}).
The game then ended with white thinking that Qe8 would be checkmate, but black can block with the bishop, and so white is left with no attack and material loss, resigning (Figure \ref{fig:game1_end}).
Overall, an aggressive opening, many tactics, blunders and checkmating attacks, this game should definitely be categorized as ``aggressive''.

\begin{figure}[H]
\centering
\begin{minipage}{0.31\textwidth}
    \centering
    \includegraphics[width=\linewidth]{game1_opening.pdf}
    \caption{After 7 moves:\\Equal position}
    \label{fig:game1_opening}
\end{minipage}
\begin{minipage}{0.31\textwidth}
    \centering
    \includegraphics[width=\linewidth]{game1_middle.pdf}
    \caption{After 10 moves:\\Black is better}
    \label{fig:game1_middle}
\end{minipage}
\begin{minipage}{0.31\textwidth}
    \centering
    \includegraphics[width=\linewidth]{game1_end.pdf}
    \caption{After 16 moves:\\Black is winning}
    \label{fig:game1_end}
\end{minipage}
\end{figure}

The second game\footnote{\url{https://lichess.org/d3UVZ3Th\#7}} had an extremely passive opening from both players (Figure \ref{fig:game2_opening}), then turned into a very closed position, with all pawns on the board and white trying to take space (Figure \ref{fig:game2_middle}).
The game ends with a simple one-move blunder by white (Figure \ref{fig:game2_end}) and an unexpected loss for black due to running out of time (which however is not relevant to our analysis).
All in all, a positional game, with very few trades and more piece maneuvering, which should definitely be part of the second category.

\begin{figure}[H]
\centering
\begin{minipage}{0.32\textwidth}
    \centering
    \includegraphics[width=\linewidth]{game2_opening.pdf}
    \caption{After 4 moves:\\Equal position}
    \label{fig:game2_opening}
\end{minipage}
\begin{minipage}{0.32\textwidth}
    \centering
    \includegraphics[width=\linewidth]{game2_middle.pdf}
    \caption{After 17 moves:\\Black is better}
    \label{fig:game2_middle}
\end{minipage}
\begin{minipage}{0.32\textwidth}
    \centering
    \includegraphics[width=\linewidth]{game2_end.pdf}
    \caption{After 20 moves:\\Black is winning}
    \label{fig:game2_end}
\end{minipage}
\end{figure}

To quantitatively define the categories, we used a set of 10 features, which can be efficiently extracted from the way the dataset is stored.
We used the monthly database of games from Lichess.org\footnote{\url{https://database.lichess.org/}}, more specifically the games between January and June 2020.
These games were part of the training data for the Maia and Maia2 engines, and the authors also provide the games from 2020 as .csv files\footnote{\url{http://csslab.cs.toronto.edu/datasets/\#maia_kdd}}.
The header of these files is shown in Table \ref{tab:lichess_csv_header}, and the information we used to calculate the features is highlighted in bold.

\begin{table}[H]
\hspace*{-1cm}
\begin{threeparttable}
\caption{CSV Header Fields for Chess Game Dataset}
\label{tab:lichess_csv_header}
\begin{tabular}{|l|l|l|l|l|}
\toprule
game\_id & type & result & white\_player & black\_player \\ \hline
white\_elo & black\_elo & time\_control & \textbf{num\_ply} & termination \\ \hline
white\_won & black\_won & no\_winner & move\_ply & move \\ \hline
cp & \textbf{cp\_rel}\* & cp\_loss & \textbf{is\_blunder\_cp} & winrate \\ \hline
winrate\_elo & winrate\_loss & is\_blunder\_wr & opp\_winrate & white\_active \\ \hline
active\_elo & opponent\_elo & active\_won & \textbf{is\_capture} & clock \\ \hline
opp\_clock & clock\_percent & opp\_clock\_percent & \textbf{low\_time}\*\* & board \\ \hline
\textbf{active\_bishop\_count} & \textbf{active\_knight\_count} & \textbf{active\_pawn\_count} & \textbf{active\_queen\_count} & active\_rook\_count \\ \hline
\textbf{is\_check} & \textbf{num\_legal\_moves} & opp\_bishop\_count & opp\_knight\_count & opp\_pawn\_count \\ \hline
opp\_queen\_count & opp\_rook\_count &  &  &  \\
\bottomrule
\end{tabular}
\begin{tablenotes}
\small
\item \* ``cp\_rel'' means the Stockfish engine evaluation relative to the side to move, measured in centipawn
\item \*\* While Maia2 does not use moves played in low time, we did, because in low time people play the most natural moves
\end{tablenotes}
\end{threeparttable}
\end{table}

So as not to have to parse the PGN of each game to extract features, which would have been very time-consuming, we only used features that can be extracted directly from the .csv files (Table \ref{tab:feature_table}).
Behind each feature there is a clear intuition as to how it should impact the categorization of the game, which is written in the third column.
For example, a longer game (with more moves) is much more susceptible to being more positional, which is why a higher num\_ply should steer the game more towards the ``positional'' category.
The objective was to calculate for each game an aggression score as the weighted sum of these features.

\begin{table}[H]
\centering
\begin{threeparttable}
\caption{Features from CSV and their Intuitive Meaning}
\label{tab:feature_table}
\begin{tabular}{lll}
\toprule
\textbf{Feature} & \textbf{Definition} & \textbf{Higher = ?} \\
\midrule
captures\_per\_move & $\frac{\texttt{captures}}{\texttt{num\_moves}}$ & tactical\\
checks\_per\_move & $\frac{\texttt{checks}}{\texttt{num\_moves}}$ & tactical\\
vol\_cp & Std.\ dev.\ of \texttt{cp\_rel} & tactical\\
blunder\_rate\_cp & $\frac{\texttt{blunders\_cp}}{\texttt{num\_moves}}$ & tactical\\
low\_time\_frac & $\frac{\texttt{low\_time\_moves}}{\max(\texttt{num\_moves},1)}$ & tactical\\
avg\_legal\_moves & Mean of \texttt{num\_legal\_moves} & positional\\
num\_ply & \texttt{num\_ply} & positional\\
queen\_activity & Mean of \texttt{active\_queen\_count} & tactical \\
minor\_piece\_loss & sum of $|\Delta(\texttt{active\_bishop\_count} + \texttt{active\_knight\_count})|$ & tactical\\
pawn\_structure\_vol & Std.\ dev.\ of \texttt{active\_pawn\_count} & tactical\\
\bottomrule
\end{tabular}
\begin{tablenotes}
\small
\item ``captures'', ``checks'', etc.\ resemble the number of lines where ``is\_capture'', ``is\_check'', etc.\ is True
\end{tablenotes}
\end{threeparttable}
\end{table}

We applied K-Means clustering on the whole dataset, recording the weights of each feature both for the centroids and for the inter-cluster boundaries (Table \ref{tab:weights_table}).
If all weights correspond to their intuitive meaning, the features that have ``tactical'' in the third column in Table \ref{tab:feature_table} should be positive, and the others negative.
There were two features for which this was not the case, namely ``queen\_activity'' and ``minor\_piece\_loss''.
The most likely explanation is queen endgames, which are extremely long and involve lots of queen moves, and games with a lot of trades that fizzle out into endgames, which even though are positional, have high minor piece trades.
Most importantly, the values of both of these features' weights were small, so even though they were on the unexpected side of the spectrum, their influence on the final aggression score was overall small.
Instead, the greatest influence had ``blunder\_rate\_cp'', ``checks\_per\_move'', and ``captures\_per\_move'', which was the expected result.

\begin{table}[H]
\centering
\caption{Features and their Weights on Aggression Score (positive steer towards aggressive, negative towards positional)}
\label{tab:weights_table}
\begin{tabular}{c|c}
\hline
\textbf{Feature} & \textbf{Weight} \\
\hline
captures\_per\_move & 0.466\\
checks\_per\_move & 0.550\\
vol\_cp & 0.369\\
blunder\_rate\_cp & 0.747\\
low\_time\_frac & 0.178\\
avg\_legal\_moves & 0.030\\
num\_ply & -0.062\\
queen\_activity & 0.155\\
minor\_piece\_loss & -0.163\\
pawn\_structure\_vol & 0.113\\
\hline
\end{tabular}
\end{table}

Then, to visualize all this, we applied a 3-dimensional principal component analysis (PCA) (Figure \ref{fig:PCA}), which would also say how well the 10 features describe the dataset.
The reason why we chose a 3D instead of the more typical 2D PCA was first, that this way more variance would be described and second, because the data is nicely separated in all dimensions.
We did this separately for each of the six months, firstly because there is more than enough data for the clustering to work well (over 2 million games), and secondly to see whether results would significantly differ between any of the months.
Because the number of games is similar in each month, statistically the results should be very similar, so a difference could point out possible problems in our thought process.
Fortunately this was not the case, and for this reason we only show the results for January 2020.
\begin{figure}[H]
%\hspace{-3cm}
\includegraphics[width=\textwidth]{PCA2.pdf}
\caption{3D PCA January 2020}
\label{fig:PCA}
\end{figure}
As can be seen, the data is very well separated in all three dimensions, which shows that the selected features were indeed relevant for describing the games.
The clusters were also very well defined, with little overlap between them, which proves that the categorization was effective.
They were also quite balanced, with the mixed cluster being the smallest, which was expected.

The next step was to train the models on the different datasets.
For this, we directly employed the training pipeline from Maia2 \cite{maia2}, only changing the datasets used for training.
For each month, we created three datasets, one for each style of play.
In order to have balanced datasets and to avoid biasing the models too much towards one style, we combined the games clustered as aggressive with those clustered as mixed to create the aggressive dataset, and similarly for the defensive dataset.
Then the mixed dataset contained all the games from the month, which also served as a baseline to compare against the official Maia2 model as well.
We used the same hyperparameters as in Maia2, training each model for only 1 epoch, as the authors found that this was sufficient for convergence.
Maia2 uses four loss functions, one of which (Total Loss) is the weighted sum of the other three.
The development of the loss functions during the training of the general model (including all the games) is shown in Figure \ref{fig:loss_stylometry}.
It is the most relevant model to show, as it used the most data and thus, the convergence can be best observed.
The Maia loss is the loss of the Maia head, which is the main head of the network, and it is the one that is used for move prediction.
It drops quickly in the first few steps, then continues to drop more slowly until about 1.8 million games, from where it only fluctuates around a certain value, which shows that the model has converged.
This shows that Maia2 learns very quickly, with the high amount of data only helping to refine the model, but not being necessary for it to learn the general patterns of human play.
This plays a greater role in the adversarial training method and is discussed in more detail in Chapter \ref{chapter:adversarial_training}.

\begin{figure}[H]
\centering
\includegraphics[width=\textwidth]{loss_stylometry2.pdf}
\caption{Development of loss functions during training of the general model. Losses quickly drop, converging after less than 3 million games.}
\label{fig:loss_stylometry}
\end{figure}

\section{Results}
\label{subsec:results_behavior_stylometry}
Evaluating human-aligned chess engines is a difficult task, as there is no clear metric to measure how well the engine mimics human play.
The typical approach is simply counting the number of times the engine's top choice matches the human move, thus measuring accuracy.
However, there are many limitations to this approach, as it does not take into account the quality of the move, nor does it consider the context of the game.
There are positions where multiple moves are equally good, or where the human move is a blunder, and in these cases, simply counting matches does not provide a complete picture of the engine's performance.
Accuracy is still a useful metric though, as it provides a baseline for comparison between different models.
Thus, we evaluated the trained models on the test set from Maia2 \cite{maia2}, comparing them with the official Maia2 model, which anyone can directly import and use.
There are three results to report here: the accuracy of each model on the whole test set (Figure \ref{fig:overall_accuracy}), the accuracy in the opening (first 5 moves) (Figure \ref{fig:opening_accuracy}), and the accuracy in the endgame (after 60 moves) (Figure \ref{fig:endgame_accuracy}).
All accuracies are reported as the percentage of times the engine's top choice matches the human move, only considering games with opponents in the same Elo bin as the engine (<100 Elo difference).
This can prove slightly misleading, as all models output a probability distribution over multiple moves, and thus even if the top choice does not match the human move, the human move might still be among the top choices of the engine.
\begin{figure}[H]
\centering
\includegraphics[width=\textwidth]{Overall_Accuracy.pdf}
\caption{Overall Accuracy. The mixed model performs best, even outperforming the official model despite being trained on less data. Tactical performs better at lower levels, while positional performs better at higher levels.}
\label{fig:overall_accuracy}
\end{figure}
\begin{figure}[H]
\centering
\includegraphics[width=\textwidth]{Opening_Accuracy.pdf}
\caption{Opening Accuracy of Models. All three behavior stylometry models significantly outperform the official Maia2 model, showing that training on opening data is indeed useful for improving human-alignment in the opening.}
\label{fig:opening_accuracy}
\end{figure}
\begin{figure}[H]
\centering
\includegraphics[width=\textwidth]{Endgame_Accuracy.pdf}
\caption{Endgame Accuracy of Models. The official Maia2 model performs best, likely due to the much larger training dataset.}
\label{fig:endgame_accuracy}
\end{figure}
As can be seen, out of our three models, the mixed one performs best, which was to be expected, as it was trained on the most data, and it even outperforms the official Maia2 model at almost all levels.
Interestingly, all three models outperform the official Maia2 model overall for levels under 1400 Elo, which shows that behavior stylometry can indeed improve human-aligned chess engines.
The choice to train on opening data as well seems to have paid off, as the opening accuracy is significantly higher for all three models compared to the official Maia2 model.
The endgame is where the official Maia2 model performs best, though only slightly better than the behavior stylometry models, which is likely due to having seen more endgame positions during training, as the behavior stylometry models had less data to learn from.
It is also interesting to note that the tactical model performs better at lower levels, while the defensive model performs better at advanced games, which aligns with the intuition that lower-level games have more tactics and blunders, while higher-level games are more positional.
Overall, these quantitative results are promising, showing that behavior stylometry can indeed improve human-aligned chess engines, but as previously mentioned, accuracy is not a perfect metric, and further research is needed to develop more comprehensive evaluation methods.

\clearpage
To better understand the differences between the models, we also analyzed their behavior in specific positions.
Here is a very descriptive example, in a quite popular position that arises from the King's Indian Defense (Figure \ref{fig:example_position}).
\begin{figure}[H]
\centering
\includegraphics[width=0.5\textwidth]{Comparison1.pdf}
\caption{Example Position and Model Predictions. Positional model chooses the most popular move Nf3, aggressive model suggests Qd2 looking to weaken black's king, official and mixed model predict Bd3}
\label{fig:example_position}
\end{figure}

The official Maia2 model and the mixed model both choose Bd3, which is the second most played move in the database, and a solid, flexible developing move.
The defensive model chooses the most played move, both at the amateur and master level, Nf3, which is also Stockfish's top choice.
It is a solid move, preparing to castle, close the center and slowly make progress on the queenside, all typical ideas in the King's Indian Defense, though definitely one of the safer options.
The aggressive model chooses the third most popular move, Qd2, with the idea of playing Bh6 and trading off the bishop, which is a key defender of black's king.
The move is definitely more ambitious, aiming to follow up with moves like g4, creating a kingside attack.

All in all, none of the engines do anything radical, but they each suggest slightly different approaches, reflecting the desired playing styles.
This meets the original expectations quite well, proving behavior stylometry to be a viable method in improving human-aligned chess engines.
While the results are promising, further research and refinement are needed to fully realize the potential of this approach, in particular regarding the integration of more nuanced behavioral traits and the evaluation metrics used to assess human-likeness in chess engines.
\chapter{Adversarial Training}
\label{chapter:adversarial_training}

\section{Related Work}
\label{section:related_work_adversarial_training}
Generative adversarial networks (GANs) were introduced by Goodfellow et al. in 2014 \cite{GAN_Goodfellow} as a framework to train generative models.
The core idea behind GANs is to pit two neural networks (a generator and a discriminator) against each other, where the generator tries to create realistic data samples, while the discriminator attempts to distinguish between real and generated samples.
In this way, both networks improve each other's performance.
Since their introduction, GANs have been widely adopted in various domains, including computer vision, healthcare, finance or natural language processing \cite{GAN_uses}.

In the domain of chess-playing AI agents, there is very limited research on the application of GANs.
The work that does exist primarily focused on using GANs to generate realistic chess positions \cite{GAN_positions} or to recognize a board position from an image \cite{GAN_image}.
There are two pieces of research that we have been able to find that explored the use of GANs in training chess-playing agents.
\cite{GAN_style} focused, similarly to some of the related work for behavior stylometry (Section \ref{section:related_work_behavior_stylometry}), on emulating the playing style of a specific human player, which, while not completely applicable in our case, did provide useful insights into the structure of the GAN and training procedure.
\cite{GAN_anticheat} explored the use of GANs in detecting cheating.
Thus, unlike in our work, the main focus of the paper was the discriminator, which was trained to distinguish between human and engine-generated moves, while the generator's role was to play moves as close to Stockfish as possible, while appearing human-like.
The results presented in the paper indicate that adversarial training can lead to more human-like play, and the conditions under which the training is performed can significantly impact the quality of the generated moves.
The discriminator, for example, had to be pre-trained, otherwise the generator failed to learn effectively.
Also, the introduction of pure Stockfish-generated moves into the training data for the generator lead to improved performance, as it provided a stronger signal for the generator to learn from.

The main insights relevant to our work were the theoretical framework for using GANs to train chess-playing agents, in particular the WGAN architecture (used in \cite{GAN_style}), and the need for changing the output of the generator from an argmax of the move probabilities to some sort of a softmax, which was necessary to allow the gradients to flow back from the discriminator to the generator.

\section{Methodology}
\label{section:methodology_adversarial_training}
The main idea behind our approach was to use Maia2 as the generator and to train a discriminator to classify moves as either human or engine-generated.
Inherently, this starting point implied an STGAN model (as introduced in \cite{GAN_style}), meaning that the generator loss was not structured purely in terms of the discriminator.
Instead, the generator already had a task-specific loss (the cross-entropy loss for move prediction that Maia2 already optimizes for), to which we added an adversarial loss component that encourages the generator to produce moves that are classified as human by the discriminator.
Combining this with the fact that we are dealing with probability distributions over moves, we chose to use the same WGAN architecture as in \cite{GAN_style}, which was designed to provide a more stable training process for GANs, and was particularly well-suited for cases where the generator's output is a probability distribution \cite{WGAN}.

Maia2's architecture relies on a ResNet backbone, which processes the input chess position and produces a feature representation that is then passed to the Elo-aware attention mechanism, which produces the final move probabilities.
The discriminator's architecture (Figure \ref{fig:discriminator_architecture}) used the same ResNet input processing (in order to keep a clear consistency between the two networks and their inputs), but then only took the move embedding and passed it through a cross-attention mechanism with the position features, followed by a final classification head that would output a linear critic score representing the probability of the move being human-generated.
Because the discriminator is essentially a classifier (but with negative values for human moves and positive ones for engine-generated moves), it minimizes the Wasserstein distance between the distribution of human moves and the distribution of generated moves through the critic loss:
\begin{equation}
\begin{aligned}
    \mathcal{L}_D = &- \mathbb{E}_{(board, human\_move) \sim P_{data}} [D(board, human\_move)]\\ 
    &+ \mathbb{E}_{(board, G\_move) \sim P_g} [D(board, G\_move)]
\end{aligned}
\end{equation}
where $P_{data}$ is the distribution of human moves and $P_g$ is the distribution of generated moves.
\begin{figure}[H]
    \hspace{-2cm}
    \includegraphics[width=1.3\textwidth]{Discriminator_Architecture.pdf}
    \caption{Architecture of the discriminator: same board processing as Maia2 flowing into a move embedding instead of skill embedding, followed by a cross-attention mechanism and a final classification head.}
    \label{fig:discriminator_architecture}
\end{figure}

Maia2's original loss is the weighted sum of the cross-entropy loss for move prediction, the BCE loss on side information and the MSE loss for the win probability (value head).
To this, we added the adversarial loss component, which encourages the generator to produce moves that are classified as human by the discriminator, thus creating the adversarial generator loss:
\begin{equation}
    \mathcal{L}_G = \mathcal{L}_{original} - \lambda \cdot \mathbb{E}_{(board, G\_move) \sim P_g} [D(board, G\_move)]
\end{equation}
where $\mathcal{L}_{original}$ is the original loss used to train Maia2, and $\lambda$ is a hyperparameter that controls the strength of the adversarial loss component.
After experimentation with different $\lambda$ values, we reached the conclusion that $\lambda = 1$ was a good choice for our training process, with $\mathbb{E}_{(board, G\_move) \sim P_g} [D(board, G\_move)]$ not reaching absolute values higher than 3 during training.

An interesting problem arose from the fact that the generator's output is a probability distribution over moves, while the discriminator only takes a single move as input.
The obvious solution was to simply use the move with the highest probability as the generator's move, but the argmax function is not differentiable, which means that the gradients from the discriminator would not be able to flow back to the generator during training.
To solve this problem, some sort of gradient estimator is needed.
We used the Gumbel-Softmax trick \cite{Gumbel_Softmax}, which allowed us to sample from a categorical distribution in a differentiable way.

The full training process is visualized in Figure \ref{fig:adversarial_training_process}.
\begin{figure}[H]
    \centering
    \includegraphics[width=\textwidth]{single_model_losses.pdf}
    \caption{Adversarial training process: stable, no mode collapse or overpowering discriminator. Generator converges slower than original Maia2 because the discriminator's signal only goes to 0 after more than 2 million games.}
    \label{fig:adversarial_training_process}
\end{figure}

The interesting part to note is that the discriminator's penalty loss is increasing during the phase where the generator's loss is drastically decreasing, and then both losses stabilize.
This means that the defining training phase is until about one million games, after which the generator has learned to produce moves that are classified as human by the discriminator, and the discriminator is no longer able to effectively distinguish between human and generated moves, thus leading to a stabilization of both losses.
This is a good sign that the adversarial training is working as intended, without typical GAN problems such as mode collapse (where the generator produces a limited variety of outputs) or the discriminator overpowering the generator (where the discriminator becomes too good at distinguishing between real and generated samples, leading to vanishing gradients for the generator).
So in terms of training stability and convergence, the adversarial training process seems to have been successful, with both the generator and discriminator improving together until they reach a point where neither can significantly outperform the other, which is the ideal outcome in GAN training.

\section{Results}
\label{section:results_adversarial_training}
Just like in the case of behavior stylometry, we evaluated the performance of the adversarially trained model using accuracy as quantitative metric, and also by looking at some qualitative examples of the generated moves.
Figure \ref{fig:top1_accuracy} shows the accuracy of the adversarially trained model compared to the original Maia2, taking into account the top predicted move.
\begin{figure}[H]
    \centering
    \includegraphics[width=\textwidth]{Top1_accuracy.pdf}
    \caption{Top-1 accuracy of the adversarially trained model compared to the original Maia2: similar accuracy at lower levels, but significantly worse performance at higher Elo.}
    \label{fig:top1_accuracy}
\end{figure}
The results indicate that the accuracy is similar between the two models at lower Elo, but at higher levels, the adversarially trained model performs worse than the original Maia2, with the gap increasing as the Elo increases.
Qualitative analysis quickly revealed the reason: the GAN model is very prone to producing blunders (moves that lead to a significant decrease in the position evaluation, such as losing material or getting checkmated).
Blunders are very common at low levels, which is why the accuracy is similar between the two models at low Elos, but they become increasingly rare at higher Elos, which is why the GAN model's performance drops significantly at higher Elos.
Figure \ref{fig:blunder} shows an example of a blunder produced by the GAN model.
\begin{figure}[H]
    \centering
    \includegraphics[width=0.5\textwidth]{Blunder-prone.pdf}
    \caption{GAN predicts Qxc6, which is then capturable by the bishop.}
    \label{fig:blunder}
\end{figure}
While it is difficult to argue that Qxc6 is not a human-like move, it is also true that it is a very bad move, which is not something we would expect from a human player at a high Elo.
The problem is that this was not an isolated case, but rather a common occurrence in the GAN model's predictions.
As far as we could tell, it represents the only notable difference between the two models, which are otherwise very similar in the moves they produce.
This is seen in the accuracy metrics as well, where the top-1 accuracies differ much more than the top-3 (Figure \ref{fig:top3_accuracy}), which means that even when the GAN model produces a blunder as its top move, it often has the correct move in its top-3 predictions, which indicates that the model is still learning to produce human-like moves, but it is not effectively learning to avoid blunders.
\begin{figure}[H]
    \centering
    \includegraphics[width=\textwidth]{Top3_accuracy.pdf}
    \caption{Top-3 accuracy of the adversarially trained model compared to the original Maia2: closer to each other than the top-1 accuracy.}
    \label{fig:top3_accuracy}
\end{figure}

The most natural explanation for this phenomenon is that the discriminator learned to classify objectively good moves as engine-generated and blunders as human, then rewarding the generator for predicting blunders.
Interestingly however, when inputting Qxc6 into the discriminator, it outputs with probability 0.986 that the move is engine-generated, which is the contrary to the hypothesis that the discriminator is rewarding blunders.
The true explanation could only be seen by retracing the training process.
When inputting Qxc6 into the discriminator produced after just 5000 games of training, it outputs with probability 0.984 that the move is human.
So while the discriminator did learn throughout the training process to distinguish much better between human and engine-generated moves, it was too late for the generator, which had already learned to produce blunders that are classified as human by the discriminator in the early stages of training, and was not able to unlearn this behavior later on, when the discriminator became more accurate.
This is a common problem in GAN training, and there are various techniques used to mitigate it.
The most common one is to pre-train the discriminator before starting the adversarial training, which we have tried and will be discussed in the next section.
Another frequently employed technique is to use a more complex architecture for the discriminator, but the fact that the discriminator after training classifies blunders correctly as engine-generated, and the fact that the training process seems to be stable and converging, indicates that the architecture is not the main problem in our case.

There are also examples when the GAN model produces better results than the original Maia2, particularly because of the discriminator's influence.
A very descriptive instance is in the position we first used to illustrate the difference between a human and a computer move, shown again in Figure \ref{fig:human_vs_engine2}.
\begin{figure}[H]
    \centering
    \includegraphics[width=0.5\textwidth]{Human_move2.pdf}
    \caption{Same as Figure \ref{fig:human_vs_engine}. Stockfish suggests b5. New model also suggest Kg8, defending the f7 pawn, which is a more human-like move than Nf8 suggested by the original Maia2.}
    \label{fig:human_vs_engine2}
\end{figure}
First and foremost, the discriminator classifies b5 as an engine-generated move with a probability of 0.999 which is what we want and expect.
Both the original Maia2 and the GAN model predict e5 and Rg8 as their top two moves, but the interesting part is the third recommendation.
Maia2 gives Nf8 with 0.17 probability, while the GAN model suggest Kg8, having the highest critic score (discriminator output) among the three moves.
Kg8 is certainly a more human-like move than Nf8, protecting the f7 pawn.

\subsection*{Pre-training the discriminator}
\label{section:pretraining_discriminator}
As previously mentioned, the most common solution to the problem of the generator learning to produce blunders in the early stages of training is to pre-train the discriminator, so that it can provide a stronger signal for the generator to learn from, and thus avoid learning to produce blunders.
We tried this approach, using a small dataset of the first 5000 games played in the month of July 2020 (none of which were used in the training process of the GAN model).
To generate the engine moves for the discriminator's training, we used a weaker version of Stockfish, tuned to play around the level of the corresponding player in each position.
These versions of Stockfish are available in the Python chess library \cite{python_stockfish}, and they are designed to play at a specific Elo level by limiting the depth of the search and the amount of time spent on each move.
While the moves generated by these Stockfish versions are far from perfect (many times they are artificial blunders or too strong for the corresponding Elo level), they still provide clear patterns for the discriminator to learn from.
Considering that the purpose of this pre-training was neither to achieve a very high accuracy for the discriminator, nor to create an anticheating model, accuracy was not the main metric we were looking at, but rather the loss and the critic scores for the human and engine-generated moves.

We trained the discriminator for 50 epochs, and the most interesting aspect of the pre-training was the great impact that preprocessing the dataset had on the training process.
When using the raw dataset (including opening moves, blunders, moves played with little time left on the clock etc.), the discriminator was highly unstable and unbalanced.
Critic scores for engine-generated moves were very high, and the accuracy was skewed strongly towards classifying moves as engine-generated.
This was exactly the behavior that shaped the generator's learning process in the early stages of training, leading to the blunder-prone behavior that we observed in the final model.
Cleaning the dataset by not taking opening moves (chosen as the first 10 ply) into consideration, removing blunders (defined as moves that change the Stockfish position evaluation by more than 2 centipawns) and removing moves played with less than 10 seconds left on the clock, led to a much more stable and balanced training process for the discriminator.
Interestingly, this is not instantly noticeable in the loss, represented in Figure \ref{fig:discriminator_loss} as the Wasserstein distance between the distribution of human moves and the distribution of generated moves, which is only slightly smaller for the cleaned dataset.
\begin{figure}[H]
    \centering
    \includegraphics[width=\textwidth]{Discriminator_loss.pdf}
    \caption{Wasserstein distance (inverse of discriminator loss) with and without cleaning the dataset. The difference is not significant, but the cleaned dataset does lead to a slightly smaller distance, which indicated better performance of the discriminator.}
    \label{fig:discriminator_loss}
\end{figure}

Instead, a significant difference can be observed in the accuracy of the discriminator, which is not just higher overall for the cleaned dataset (Figure \ref{fig:discriminator_total_accuracy}), but also better balanced between the two classes (Figure \ref{fig:discriminator_class_accuracy}), which is what we were looking for in order to provide a stronger signal for the generator to learn from, and thus avoid learning to produce blunders in the early stages of training.
\begin{figure}[H]
    \centering
    \includegraphics[width=\textwidth]{Discriminator_total_accuracy.pdf}
    \caption{Total accuracy of the discriminator with and without cleaning the dataset: higher accuracy for the cleaned dataset across all levels.}
    \label{fig:discriminator_total_accuracy}
\end{figure}
\begin{figure}[H]
    \centering
    \includegraphics[width=\textwidth]{Discriminator_class_accuracy.pdf}
    \caption{Class accuracy of the discriminator with and without cleaning the dataset: much more balanced predictions for the cleaned dataset.}
    \label{fig:discriminator_class_accuracy}
\end{figure}

The stability can also be observed in the critic scores the discriminator outputs.
Taking as example the blunder originally recommended by the GAN model (Qxc6 in the position shown in Figure \ref{fig:blunder}), the discriminator trained on the raw dataset outputs a critic score of almost 14, whereas the cleaned dataset leads to a much more reasonable score of 3.1.
So while both models classify the move as clearly human (which is what we want, as the move is indeed human-like), the model trained on the raw dataset is rewarding the generator much more for producing this blunder, which is what leads to the blunder-prone behavior of the GAN model.

The training process (Figure \ref{fig:training_process_cleaned}), as expected, was not much different from the one without pre-training.
The main difference is that at the very beginning of the training process, the discriminator already had a better understanding of the difference between human and engine-generated moves, which provided a stronger signal for the generator to learn from.
Thus the discriminator's signal was strongest at the start and then quickly converged towards 0 as the generator learns to produce human-like moves.
This also meant that the generator's learning was slightly accelerated at the very beginning, though the overall training process was very similar to the one without pre-training.
\begin{figure}[H]
    \centering
    \includegraphics[width=\textwidth]{GAN_clean_training.pdf}
    \caption{Adversarial training process with pre-training the discriminator on the cleaned dataset: Discriminator's signal is stronger at the start, leading to a slightly faster convergence of the generator.}
    \label{fig:training_process_cleaned}
\end{figure}

The results are also quite similar to the ones without pre-training.
Still, the only main difference to the original Maia2, from a qualitative standpoint, is the slightly blunder-prone behavior, which is still present in the GAN model with pre-training, though to a much lesser extent.
This is also reflected in the accuracy metrics, where the top-1 accuracy is still lower than the original Maia2 at higher Elos, but the gap is smaller than in the case without pre-training (Figure \ref{fig:top1_accuracy_cleaned}).
\begin{figure}[H]
    \centering
    \includegraphics[width=\textwidth]{GAN_clean_accuracy.pdf}
    \caption{Top-1 accuracy of the adversarially trained model with pre-training the discriminator on the cleaned dataset, compared to the original Maia2. The accuracy is still lower at higher Elos, but the gap is smaller than in the case without pre-training. At Elo<1100 (>50\% of players), it even outperforms the original Maia2.}
    \label{fig:top1_accuracy_cleaned}
\end{figure}

Overall, while pre-training the discriminator on a cleaned dataset is a step in the right direction, it does not completely solve the problem of the generator learning to produce blunders in the early stages of training, which indicates that there are other factors at play as well.
Intuitively, it is easy to explain this behavior from the construction of the discriminator.
Distinguishing between human and engine-generated moves is a very complex task, with a great deal of research still needed to be able to truly automate anti-cheating types of models.
Even the most modern anti-cheating mechanisms still rely nowadays on human experts' analysis, and automated models need much more data than just one move (hundreds of games played by the same player, time management etc.) to be able to achieve a good accuracy (see for example \cite{chess_anticheat}, where multiple models and neural networks are combined).
Thus, it is not surprising that the discriminator is not able to provide a strong signal for the generator to learn from, especially since Maia2 is very strong even with little training data.
While the adversarial training method has many possible applications (such as style-transfer) and advantages (better explainability, flexibility etc.), we believe it is more of a fine-tuning method to slightly modify the behavior of an already well-trained model, rather than a method that can lead to significant improvements over the original model, especially if the original model is already very strong.
Just like with behavior stylometry, we advocate for psychologically-informed metrics and methods, which could lead to much more refined discriminators and thus better results for the generator, rather than just relying on pure machine learning techniques.
\chapter{Conclusion}
\label{chapter:conclusion}

This thesis aimed to explore the potential, challenges and limitations of behavior stylometry and adversarial training for improving state-of-the-art human-aligned chess engines.

Behavior stylometry showed promising results in creating more nuanced engines by filtering the training data to only include games played in a certain style.
While our approach only used two categories (tactical and positional), we advocate for a more fine-grained categorization of chess styles in future research, as well as the use of more psychologically-grounded evaluation metrics to better capture the nuances of human play.
Generally, the results indicate that while behavior stylometry does not lead to a significant improvement in pure accuracy, it does modify the model's predictions in a way that is consistent with the intended style.
Benefits of this method include its relatively simple implementation, its stability during training and the fact that it does not require any modification to the training process, as all the work is done in the data preprocessing stage.
Potential applications could include the creation of more personalized chess engines, allowing players to better learn from, or training against, a specific style of play.

Adversarial training's main challenge is the wide range of possibilities regarding the implementation (total model design, discriminator architecture, hyperparameters etc.), which makes it difficult to fully generalize the results and draw conclusions about the method as a whole.
In our case, the WGAN setup proved to lead to a stable training process, creating a discriminator that can distinguish between human and engine-generated moves with a relatively high accuracy, but a generator that does not significantly change compared to the original Maia2 model, except for the more blunder-prone behavior.
We identified the cause of this behavior to be the fact that the discriminator is not able to capture subtle aspects of human play (such as strategic planning, long-term thinking or immediate tactical awareness), and thus focuses on more superficial features, such as blunders, which are more common in human play.
More research is needed to find ways to make the discriminator capture these more subtle aspects, for example by incorporating cognitive insights into its design, or by focusing on emulating the human thought process rather than just the final move choice.
The clear advantage of adversarial training is its explainability, as the discriminator's output adds a layer of interpretability to the model's predictions, and even becomes a tool by itself, as it can be used to evaluate the human-likeness of any move.
Additionally, the flexibility of the method allows for tuning the model to achieve different goals than what we tried, for example combining the adversarial training with behavior stylometry and skewing the generator towards a specific playing style, or making the discriminator focus on specific aspects of human play, such as tactical patterns or endgame play, thus creating a more specialized engine.

Overall, while both methods have their own advantages and limitations, we believe that they are both promising directions for future research in the field of human-aligned engines.
However, by far the greatest leap forward in the field would consist in psychological studies of human chess play.
Pure machine learning techniques are unlikely to lead to massive improvements over the current state-of-the-art models, and since the main goal of human-aligned engines is not pure accuracy, but rather personalization, explainability and flexibility, we advocate for psychologically-informed evaluation metrics and training methods as the most promising direction for future research in the field.
\printbibliography

%\appendix
%\input{latexhints-english}

\pagestyle{empty}
\renewcommand*{\chapterpagestyle}{empty}
%\Versicherung
\end{document}
